\documentclass[11pt,a4paper]{article}

% -----------------------------------------------------------------------------
% Packages
% -----------------------------------------------------------------------------
\usepackage[margin=1in]{geometry}
\usepackage[T1]{fontenc}
\usepackage[utf8]{inputenc}
\usepackage{amsmath,amssymb,amsfonts,mathtools}
\usepackage{bm}
\usepackage{booktabs}
\usepackage{array}
\usepackage{longtable}
\usepackage{graphicx}
\usepackage{xcolor}
\usepackage{hyperref}
\usepackage{enumitem}
\usepackage{listings}
\usepackage{microtype}
\usepackage{float}

\hypersetup{
    colorlinks=true,
    linkcolor=blue!60!black,
    citecolor=blue!60!black,
    urlcolor=blue!60!black
}

\lstdefinestyle{pythonstyle}{
    language=Python,
    basicstyle=\ttfamily\small,
    keywordstyle=\color{blue!70!black},
    commentstyle=\color{green!40!black},
    stringstyle=\color{red!50!black},
    showstringspaces=false,
    breaklines=true,
    frame=single,
    rulecolor=\color{black!20},
    columns=fullflexible
}
\lstset{style=pythonstyle}

% -----------------------------------------------------------------------------
% Macros
% -----------------------------------------------------------------------------
\newcommand{\R}{\mathbb{R}}
\newcommand{\E}{\mathbb{E}}
\newcommand{\Var}{\operatorname{Var}}
\newcommand{\diag}{\operatorname{diag}}
\newcommand{\Normal}{\mathcal{N}}
\newcommand{\Unif}{\mathcal{U}}
\newcommand{\given}{\,|\,}
\newcommand{\thetaS}{\boldsymbol{\theta}}
\newcommand{\xobs}{\mathbf{x}_{\mathrm{obs}}}
\newcommand{\yobs}{\mathbf{y}_{\mathrm{obs}}}
\newcommand{\sld}{\ensuremath{\mathrm{SLD}}}
\newcommand{\rel}{\ensuremath{\mathrm{rel}}}
\newcommand{\PIT}{\mathrm{PIT}}
@
\title{A Full Integrated Report for the Oncology PK/PD and Simulation-Based Inference Curriculum\\\large Notebook-by-Notebook Theory, Outputs, and Identifiability Lessons}
\author{Technical Report}
\date{\today}

\begin{document}
\maketitle

\begin{abstract}
This report expands the original write-up for \texttt{00\_clinical\_audit\_and\_model\_discovery.ipynb} into a full technical document covering the complete notebook sequence through Notebook 06B. The project starts with heterogeneous longitudinal oncology trial files, builds a canonical tumor-size and dose layer, fits mechanistic tumor-growth inhibition models, studies their hold-out and rolling-origin forecasting behavior, calibrates predictive uncertainty, selects clinically coherent simulation-based inference (SBI) starter strata, constructs real-template synthetic simulators, and then tests neural posterior estimation (NPE) on anti-angiogenic synthetic data. The central result is not a success story about increasingly complex neural models. The central result is an identifiability diagnosis: sparse relative \sld{} trajectories can support useful empirical forecasting and modest posterior predictive uncertainty quantification, but they weakly identify detailed mechanistic PK/PD parameters even under simulator-matched synthetic data. Consequently, the correct next milestone is posterior predictive validation and uncertainty calibration, not raw parameter recovery on real patients.
\end{abstract}

\tableofcontents
\newpage

% -----------------------------------------------------------------------------
\section{High-Level Problem Statement}
% -----------------------------------------------------------------------------

We observe cancer patients indexed by
\[
    i=1,\ldots,N.
\]
Patient $i$ receives one or more anticancer drugs over time. Let the dose history be
\[
    u_i=\{(t_{ik}^{\mathrm{dose}},d_{ik},a_{ik})\}_{k=1}^{K_i},
\]
where $t_{ik}^{\mathrm{dose}}$ is the administration time, $d_{ik}$ is the amount, and $a_{ik}$ identifies the drug or regimen.

Tumor burden is observed at scan times
\[
    t_{ij}^{\mathrm{scan}}, \qquad j=1,\ldots,T_i,
\]
through an \sld-like scalar $y_{ij}>0$. The relative tumor trajectory is
\[
    r_{ij}=\frac{y_{ij}}{y_{i1}},
\]
and on log scale
\[
    \ell_{ij}=\log r_{ij}=\log y_{ij}-\log y_{i1}.
\]

The scientific target is a patient-specific latent parameter vector
\[
    \theta_i\in\Theta\subset\R^p
\]
containing quantities such as growth rates, transit rates, treatment effect sizes, resistance fractions, or carrying-capacity parameters. The eventual Bayesian inverse problem is
\[
    p(\theta_i\given x_i),
\]
where $x_i$ encodes scan times, relative tumor sizes, dose history, and covariates.

This report explains why the project did \emph{not} jump directly to SBI. The full workflow was
\[
\boxed{
\text{clinical audit}
\rightarrow
\text{mechanistic fitting}
\rightarrow
\text{hold-out validation}
\rightarrow
\text{rolling-origin baselines}
\rightarrow
\text{probabilistic calibration}
\rightarrow
\text{stratum selection}
\rightarrow
\text{synthetic SBI}
\rightarrow
\text{identifiability diagnosis}
}
\]
rather than
\[
    \text{raw longitudinal data}\rightarrow\text{neural posterior estimation immediately}.
\]

The logic is simple: if the simulator family is structurally inadequate, or if the data weakly identify its parameters, a learned posterior can look sophisticated while remaining scientifically misleading.

% -----------------------------------------------------------------------------
\section{Global Mathematical Framing}
% -----------------------------------------------------------------------------

\subsection{Observation and Loss Geometry}

Throughout the notebook sequence, the main residual is defined in log-relative tumor space:
\[
    e_j(\theta)=\log r_j-\log \widehat r(t_j;\theta,u).
\]
This is appropriate because multiplicative tumor-size errors are more clinically realistic than additive millimeter errors. A symmetric doubling/halving discrepancy gives equal-magnitude residuals with opposite sign.

The standard log-RMSE is
\[
    \operatorname{RMSE}_{\log}(\theta)
    =
    \sqrt{\frac{1}{T}\sum_{j=1}^{T}e_j(\theta)^2}.
\]
A value near $0.30$ corresponds to a typical multiplicative error of
\[
    \exp(0.30)-1\approx 35\%.
\]

\subsection{Forward-Simulator View}

Every notebook after the audit phase is implicitly organized around a forward simulator
\[
    \theta\sim p(\theta),
    \qquad
    x\sim p(x\given\theta),
\]
even when the immediate task is deterministic fitting. The SBI question only becomes meaningful once we know:
\begin{enumerate}[leftmargin=2em]
    \item what the latent dynamics are,
    \item what the observation model is,
    \item what forecasting baselines the simulator must beat,
    \item whether the parameters are even identifiable from the chosen observation channel.
\end{enumerate}

The rest of the report follows the notebook chronology while keeping this inverse-problem perspective explicit.

% -----------------------------------------------------------------------------
\section{Notebook 00: Clinical Audit and First Model Discovery}
% -----------------------------------------------------------------------------

Notebook 00 is the foundation of the whole project. It is the notebook originally documented in \texttt{notebook1.tex}, so this section preserves its depth and then connects it to the later notebooks.

\subsection{Canonical Clinical Data Layer}

The raw data tree contains multiple oncology trials with different schemas, file names, and measurement conventions. Each discovered study $s$ is mapped into a canonical object
\[
    \mathcal D_s=(O_s,D_s,C_s),
\]
where
\begin{itemize}[leftmargin=2em]
    \item $O_s$ is the canonical tumor observation table,
    \item $D_s$ is the canonical dose table,
    \item $C_s$ is the canonical covariate table.
\end{itemize}

The observation table contains columns such as
\[
\begin{array}{ll}
\texttt{study\_id}, & \texttt{patient\_id},\\
\texttt{raw\_day}, & \texttt{time\_days},\\
\texttt{sld\_mm}, & \texttt{rel\_sld},\\
\texttt{baseline\_sld\_mm}, & \texttt{n\_lesions},\\
\texttt{arm}. &
\end{array}
\]
The relative tumor burden is computed patient-wise as
\[
    \texttt{rel\_sld}_{ij}
    =
    \frac{\texttt{sld\_mm}_{ij}}
    {\texttt{baseline\_sld\_mm}_{i}}.
\]

The dose table stores
\[
    \texttt{study\_id},
    \texttt{patient\_id},
    \texttt{time\_days},
    \texttt{drug},
    \texttt{dose\_amount},
    \texttt{dose\_unit},
    \texttt{arm},
\]
while the covariate table keeps demographic and arm metadata for later hierarchical priors or covariate-conditioned SBI.

\subsection{Dataset Discovery and Eligibility}

The loader recursively searches for known trial-specific patterns and maps them into the common schema. A patient is considered PK/PD-eligible when the trajectory has enough scans and a plausible baseline near treatment initiation. Typical filters are of the form
\[
    T_i\geq 3,
    \qquad
    t_{i1}\in[-60,14]\ \text{days}.
\]

Duplicate same-day observations are collapsed, typically via a within-day median,
\[
    \widetilde y_i(t)=\operatorname{median}\{y_{ij}:t_{ij}=t\},
\]
so that the ODE solver only sees strictly increasing evaluation times.

\subsection{Audit Outputs}

Notebook 00 loaded ten studies and produced the canonical audit summarized in Table~\ref{tab:audit-summary}. The table is important because it explains why later notebooks focus on certain strata and not others: scan counts, dose completeness, and tumor-size distributions differ meaningfully across studies.

\begin{table}[H]
\centering
\small
\caption{Canonical clinical audit summary from Notebook 00.}
\label{tab:audit-summary}
\begin{tabular}{lrrrrrrr}
\toprule
Study & Obs. pts. & Eligible pts. & Obs. rows & Dose rows & Median scans & Median baseline \sld{} & Median rel.\ \sld{} \\
\midrule
A6181122 & 378 & 321 & 2,307 & 23,046 & 6.0 & 115.5 & 0.807 \\
DATASET\_263 & 945 & 671 & 4,139 & 45,516 & 4.0 & 120.0 & 0.949 \\
DATASET\_264 & 935 & 804 & 5,369 & 92,188 & 5.0 & 134.0 & 0.673 \\
EFC10262 & 592 & 446 & 2,953 & 30,117 & 4.0 & 110.3 & 0.969 \\
EFC5505 & 432 & 358 & 2,001 & 24,901 & 5.0 & 110.0 & 0.685 \\
EFC4972 & 320 & 273 & 1,503 & 23,639 & 5.0 & 90.2 & 0.737 \\
AMGEN\_20050203 & 488 & 387 & 2,476 & 0 & 4.0 & 157.5 & 0.672 \\
AMGEN\_20020408 & 275 & 69 & 672 & 0 & 2.0 & 196.0 & 1.000 \\
CEDIRANIB\_HORIZONIII & 688 & 584 & 3,340 & 46,174 & 5.0 & 10.95 & 0.759 \\
AMGEN\_20040249 & 809 & 655 & 3,802 & 67,273 & 4.0 & 127.0 & 0.794 \\
\bottomrule
\end{tabular}
\end{table}

\subsection{Exposure Approximation}

The first mechanistic models do not yet use full pharmacokinetics. Instead, dose events are converted into a simple concentration-like forcing curve
\[
    C(t)=\sum_{k:t_k\leq t}d_k\exp[-k_{\mathrm{el}}(t-t_k)].
\]
Dose amounts are normalized by the median positive dose amount so that $k_{\mathrm{kill}}$ can be given broad but stable priors. This simplification is deliberate: the project first tests whether the tumor-size signal alone supports meaningful inverse inference before adding additional PK complexity.

\subsection{Baseline Simeoni Tumor-Growth Inhibition Model}

The baseline simulator is a four-state Simeoni-style model with proliferating tumor $x_1$ and three damaged/transit compartments $x_2,x_3,x_4$:
\[
    V(t)=x_1(t)+x_2(t)+x_3(t)+x_4(t).
\]
The patient-specific parameter vector is
\[
    \theta=(\lambda_0,\lambda_1,\psi_g,k_{\mathrm{kill}},k_{\mathrm{tr}},v_0).
\]

The saturating growth function is
\[
    g(x_1;\lambda_0,\lambda_1,\psi_g)
    =
    \frac{\lambda_0 x_1}
    {\left[1+\left(\frac{\lambda_0}{\lambda_1}x_1\right)^{\psi_g}\right]^{1/\psi_g}}.
\]
For small $x_1$, the dynamics are approximately exponential, $g(x_1)\approx\lambda_0x_1$. For large $x_1$, the growth rate approaches an approximately linear regime, $g(x_1)\approx\lambda_1$.

The ODE system is
\[
\begin{aligned}
    \frac{dx_1}{dt}
    &=g(x_1)-k_{\mathrm{kill}}C(t)x_1,\\
    \frac{dx_2}{dt}
    &=k_{\mathrm{kill}}C(t)x_1-k_{\mathrm{tr}}x_2,\\
    \frac{dx_3}{dt}
    &=k_{\mathrm{tr}}(x_2-x_3),\\
    \frac{dx_4}{dt}
    &=k_{\mathrm{tr}}(x_3-x_4),
\end{aligned}
\]
with
\[
    x(0)=(v_0,0,0,0).
\]
The fitted observable is the relative trajectory
\[
    \widehat r(t_j;\theta)=\frac{V(t_j;\theta)}{V(t_1;\theta)}.
\]

\subsection{Optimization and Robust Residuals}

The project optimizes on transformed parameters $z=\log\theta$ to enforce positivity and improve numerical conditioning. Both squared-error and Huber-style robust losses are used in log-relative space. For threshold $\delta$,
\[
    \rho_\delta(e)=
    \begin{cases}
        \frac{1}{2}e^2, & |e|\leq\delta,\\
        \delta\left(|e|-\frac{1}{2}\delta\right), & |e|>\delta.
    \end{cases}
\]
The corresponding robust aggregate score is
\[
    L_{\mathrm{Huber}}(\theta)
    =
    \sqrt{2\cdot\frac{1}{T}\sum_{j=1}^{T}\rho_\delta(e_j(\theta))}.
\]

This matters because sparse clinical trajectories often contain outlying jumps or measurement inconsistencies. A model that only succeeds under pure least-squares may be too fragile for later SBI.

\subsection{Observation-Shift Extension}

Early residual plots revealed cases where the mechanistic shape looked plausible but the later observations appeared shifted upward or downward. A multiplicative observation changepoint was introduced:
\[
    \log \widehat r^{\mathrm{shift}}_j
    =
    \log \widehat r(t_j;\theta)+b\,\mathbf{1}_{j\geq k},
\]
where $k$ is a breakpoint index and $\exp(b)$ is the multiplicative shift factor. This is not a biological effect model. It is a measurement-process diagnostic layer.

\subsection{Nested Resistant-Simeoni Model}

The first two-compartment sensitive/resistant extension was not properly nested in the original Simeoni model and often degraded fits. The key modeling correction was the \emph{nested resistant-Simeoni} extension, which preserves the original sensitive damage chain and adds a resistant compartment:
\[
    x(t)=\bigl(s_1(t),s_2(t),s_3(t),s_4(t),r(t)\bigr),
\]
with total burden
\[
    V(t)=s_1+s_2+s_3+s_4+r.
\]

The parameters are
\[
    \theta_{\mathrm{NR}}
    =
    (\lambda_{0,s},\lambda_1,\psi_g,k_{\mathrm{kill},s},k_{\mathrm{tr}},v_0,f_{r0},m_r,q_r,k_{sr}).
\]
The resistant parameters are tied to the sensitive ones through
\[
    \lambda_{0,r}=m_r\lambda_{0,s},
    \qquad
    k_{\mathrm{kill},r}=q_r k_{\mathrm{kill},s}.
\]

The ODE becomes
\[
\begin{aligned}
    \frac{ds_1}{dt}&=g_s(s_1)-k_{\mathrm{kill},s}C(t)s_1-k_{sr}s_1,\\
    \frac{ds_2}{dt}&=k_{\mathrm{kill},s}C(t)s_1-k_{\mathrm{tr}}s_2,\\
    \frac{ds_3}{dt}&=k_{\mathrm{tr}}(s_2-s_3),\\
    \frac{ds_4}{dt}&=k_{\mathrm{tr}}(s_3-s_4),\\
    \frac{dr}{dt}&=g_r(r)-k_{\mathrm{kill},r}C(t)r+k_{sr}s_1.
\end{aligned}
\]
This model approximately reduces to the plain Simeoni model when
\[
    f_{r0}\approx 0,
    \qquad
    k_{sr}\approx 0.
\]

\subsection{Notebook 00 Outputs}

On the first showcased patient, DATASET\_263 / 203705005, using the patient-specific dose schedule reduced differential-evolution log-RMSE from roughly $0.273$ under a generic q2w schedule to roughly $0.202$ under the actual dose schedule. This established a key principle for the rest of the project: even simple tumor models benefit materially from correct dosing context.

For the selected retrospective 20-patient batch, the main results were:

\begin{table}[H]
\centering
\small
\caption{Notebook 00 retrospective model comparison over the 20-patient batch.}
\label{tab:n00-results}
\begin{tabular}{lrrrr}
\toprule
Model & Median RMSE & Mean RMSE & Min & Max \\
\midrule
FinalSelectedConservative & 0.151 & 0.181 & 0.036 & 0.530 \\
NestedResistantSimeoni\_plus\_posthoc\_shift & 0.151 & 0.176 & 0.035 & 0.485 \\
NestedResistantSimeoni & 0.178 & 0.200 & 0.036 & 0.500 \\
Simeoni\_plus\_joint\_observation\_shift & 0.238 & 0.271 & 0.058 & 0.612 \\
Simeoni & 0.290 & 0.290 & 0.084 & 0.612 \\
\bottomrule
\end{tabular}
\end{table}

The final conservative model counts were:
\begin{itemize}[leftmargin=2em]
    \item NestedResistantSimeoni\_plus\_posthoc\_shift: 7 patients,
    \item NestedResistantSimeoni: 6 patients,
    \item Simeoni: 6 patients,
    \item Simeoni\_plus\_joint\_observation\_shift: 1 patient.
\end{itemize}

The median RMSE improved from
\[
    0.2898 \quad \text{(plain Simeoni)}
\]
to
\[
    0.2405 \quad \text{(best simple existing model)}
\]
and then to
\[
    0.1508 \quad \text{(final conservative selected model)}.
\]

\paragraph{Interpretation.}
Notebook 00 delivered a strong retrospective modeling result and a strong warning at the same time. Richer dynamics and observation layers clearly improved fit. But this was \emph{in-sample} evidence only. The next notebook had to ask whether these apparent improvements generalized to future scans or simply overfit sparse trajectories.

% -----------------------------------------------------------------------------
\section{Notebook 01: Hold-Out Model Validation}
% -----------------------------------------------------------------------------

Notebook 01 is where the project first stress-tested the optimism created by Notebook 00. The question changed from
\[
    \text{Can the model explain an entire observed trajectory?}
\]
to
\[
    \text{Can the model trained on early scans predict later scans?}
\]

\subsection{Train/Test Split Logic}

For each of the same 20 patients, the trajectory was split chronologically. Roughly the first
\[
    75\%
\]
of scans formed the training segment and the last
\[
    25\%
\]
formed the hold-out segment. Models were fit only on the training prefix and evaluated on future scans.

This is critical because a sparse longitudinal fit can look impressive if it is allowed to see the entire future trajectory during optimization.

\subsection{Hold-Out Objective and Complexity Penalty}

Training still used log-RMSE-like objectives, but evaluation focused on future prediction error. Notebook 01 also tracked BIC-like complexity-aware comparison. In generic form,
\[
    \mathrm{BIC}=n\log\!\left(\frac{\mathrm{RSS}}{n}\right)+k\log n,
\]
where $k$ is the number of fitted parameters. The notebook compared:
\begin{itemize}[leftmargin=2em]
    \item rule-based conservative selection,
    \item BIC-based selection,
    \item per-model train and test RMSE summaries.
\end{itemize}

\subsection{Notebook 01 Outputs}

The key numerical summary is shown in Table~\ref{tab:n01-results}.

\begin{table}[H]
\centering
\small
\caption{Notebook 01 hold-out validation summary. Training fit improved with flexibility, but future test error remained high.}
\label{tab:n01-results}
\begin{tabular}{lrrrr}
\toprule
Model & Median train RMSE & Median test RMSE & Mean test RMSE & Median BIC \\
\midrule
NestedResistantSimeoni & 0.172 & 0.488 & 0.646 & -10.60 \\
Simeoni\_plus\_joint\_observation\_shift & 0.197 & 0.495 & 0.591 & -11.18 \\
NestedResistantSimeoni\_plus\_posthoc\_shift & 0.154 & 0.498 & 0.613 & -7.34 \\
Simeoni & 0.216 & 0.513 & 0.602 & -13.55 \\
\bottomrule
\end{tabular}
\end{table}

The conservative rule-based selection counts were:
\begin{itemize}[leftmargin=2em]
    \item Simeoni: 9,
    \item NestedResistantSimeoni: 6,
    \item NestedResistantSimeoni\_plus\_posthoc\_shift: 4,
    \item Simeoni\_plus\_joint\_observation\_shift: 1.
\end{itemize}

The BIC selection counts were:
\begin{itemize}[leftmargin=2em]
    \item Simeoni: 14,
    \item Simeoni\_plus\_joint\_observation\_shift: 3,
    \item NestedResistantSimeoni: 2,
    \item NestedResistantSimeoni\_plus\_posthoc\_shift: 1.
\end{itemize}

Most importantly, the median test RMSE of the rule-selected model was
\[
    0.4246,
\]
while the median test RMSE of the BIC-selected model was
\[
    0.5257.
\]

\subsection{What Notebook 01 Changed Scientifically}

Notebook 01 showed that the retrospective success of Notebook 00 did not transfer cleanly to future scans. Mechanistic flexibility lowered training error substantially, but median test RMSEs still clustered near
\[
    0.49\text{--}0.51,
\]
which is much worse than the training fit quality.

This is the first major identifiability warning:
\[
\boxed{
\text{a mechanistic model can interpolate early scans well while still failing as a forecasting model}
}
\]
because the data do not constrain the hidden parameters sharply enough.

That finding motivated Notebook 02, where the project stopped treating ``future prediction'' as a single split problem and switched to repeated rolling-origin forecasting.

% -----------------------------------------------------------------------------
\section{Notebook 02: Forecasting Baselines and Rolling-Origin Validation}
% -----------------------------------------------------------------------------

Notebook 02 reframed the problem in a more practical way. Instead of one train/test split per patient, it asked:
\[
    \text{Given scans }1,\ldots,j,\ \text{how well can we predict scan }j+1?
\]

\subsection{Rolling-Origin Forecasting Setup}

At each forecast origin $j$, the model sees a prefix
\[
    \{(t_1,r_1),\ldots,(t_j,r_j)\}
\]
and predicts the next relative tumor size. The target is therefore a one-step predictive distribution or point prediction for
\[
    r_{j+1}\given r_{1:j},u_{1:j}.
\]

The loss on each origin is based on the log error
\[
    \varepsilon_{j+1}
    =
    \log r_{j+1}-\log \widehat r_{j+1}.
\]
Notebook 02 reports
\begin{itemize}[leftmargin=2em]
    \item median absolute log error,
    \item mean absolute log error,
    \item one-step RMSE.
\end{itemize}

\subsection{Empirical Baselines}

Three empirical baselines are especially important:

\paragraph{LOCF.}
Last observation carried forward predicts
\[
    \widehat r_{j+1}=r_j.
\]
This is a persistence baseline.

\paragraph{RecentLogSlope.}
Use the most recent local slope on log scale,
\[
    \widehat{\log r}_{j+1}
    =
    \log r_j
    +
    \frac{\log r_j-\log r_{j-1}}{t_j-t_{j-1}}
    (t_{j+1}-t_j).
\]

\paragraph{LogLinear.}
Fit a least-squares line to the available prefix in log space,
\[
    \log r(t)\approx \beta_0+\beta_1 t,
\]
and extrapolate to the next scan time.

These baselines are scientifically useful because they set the minimum forecasting standard that a mechanistic simulator must beat before being trusted for SBI.

\subsection{Mechanistic Competitors}

Notebook 02 compares the empirical baselines against:
\begin{itemize}[leftmargin=2em]
    \item Simeoni,
    \item Simeoni\_plus\_joint\_observation\_shift,
    \item NestedResistantSimeoni,
    \item NestedResistantSimeoni\_plus\_train\_shift.
\end{itemize}

The ``train shift'' variant only uses information available up to the forecast origin, preventing direct look-ahead leakage.

\subsection{Notebook 02 Outputs}

The main one-step forecasting results are in Table~\ref{tab:n02-results}.

\begin{table}[H]
\centering
\small
\caption{Notebook 02 rolling-origin one-step forecasting on the 20-patient cohort.}
\label{tab:n02-results}
\begin{tabular}{lrrrr}
\toprule
Model & N origins & Median $|\log$ err$|$ & Mean $|\log$ err$|$ & RMSE \\
\midrule
LOCF & 134 & 0.046 & 0.102 & 0.185 \\
RecentLogSlope & 134 & 0.066 & 0.127 & 0.210 \\
LogLinear & 134 & 0.243 & 0.295 & 0.376 \\
NestedResistantSimeoni\_plus\_train\_shift & 134 & 0.234 & 0.313 & 0.404 \\
NestedResistantSimeoni & 134 & 0.239 & 0.342 & 0.462 \\
Simeoni\_plus\_joint\_observation\_shift & 134 & 0.312 & 0.411 & 0.542 \\
Simeoni & 134 & 0.320 & 0.442 & 0.582 \\
\bottomrule
\end{tabular}
\end{table}

Selection strategies were also compared:

\begin{table}[H]
\centering
\small
\caption{Notebook 02 rolling-origin selection strategy comparison.}
\label{tab:n02-selection}
\begin{tabular}{lrrr}
\toprule
Strategy & Count & Median $|\log$ err$|$ & RMSE \\
\midrule
oracle\_per\_origin & 134 & 0.026 & 0.105 \\
global\_best\_LOCF & 134 & 0.046 & 0.185 \\
adaptive\_past\_patient\_errors & 134 & 0.049 & 0.204 \\
\bottomrule
\end{tabular}
\end{table}

The per-patient best-model counts were:
\begin{itemize}[leftmargin=2em]
    \item LOCF: 12 patients,
    \item RecentLogSlope: 7 patients,
    \item Simeoni: 1 patient.
\end{itemize}

\subsection{Interpretation}

Notebook 02 changes the direction of the whole project. The lesson is not just that mechanistic models are imperfect. The lesson is that on sparse real clinical tumor trajectories, a simple persistence forecast is extremely hard to beat:
\[
\boxed{
\text{if a mechanistic model cannot beat LOCF in rolling-origin prediction, retrospective fit alone is not enough evidence}
}
\]

This is why Notebook 03 pivots to probabilistic forecasting and uncertainty calibration. Once point forecasting is difficult, the next scientifically meaningful question becomes: can we at least produce calibrated predictive intervals?

% -----------------------------------------------------------------------------
\section{Notebook 03: Probabilistic Forecasting and Calibration}
% -----------------------------------------------------------------------------

Notebook 03 makes a conceptual move rather than only a numerical one. It treats deterministic one-step predictions as mean structures and adds a residual uncertainty model on top.

\subsection{From Point Forecasts to Predictive Distributions}

Suppose a deterministic forecaster returns $\widehat r_{j+1}$. Notebook 03 models the next observation in log space as
\[
    \log r_{j+1}
    =
    \log \widehat r_{j+1}
    +
    \epsilon_{j+1}.
\]
Residuals from rolling-origin backtesting are then used to estimate the distribution of $\epsilon$.

Two main strategies are used.

\paragraph{Gaussian calibration.}
Assume
\[
    \epsilon\sim\Normal(\mu,\sigma^2),
\]
usually with $\mu$ close to zero after calibration. Then a predictive interval for level $1-\alpha$ is
\[
    \log \widehat r_{j+1}\pm z_{1-\alpha/2}\sigma.
\]
The Gaussian negative log likelihood is
\[
    \mathrm{NLL}
    =
    \frac{1}{2}\log(2\pi\sigma^2)
    +
    \frac{(\epsilon-\mu)^2}{2\sigma^2}.
\]

\paragraph{Empirical-quantile calibration.}
Instead of imposing Gaussian tails, estimate residual quantiles directly. If $q_{\alpha/2}$ and $q_{1-\alpha/2}$ are empirical residual quantiles, the predictive interval becomes
\[
    \bigl[
        \log \widehat r_{j+1}+q_{\alpha/2},
        \log \widehat r_{j+1}+q_{1-\alpha/2}
    \bigr].
\]
This is especially attractive when residuals are skewed or heavy tailed.

\subsection{Coverage, Width, and PIT Diagnostics}

Notebook 03 formalizes several metrics that remain central later:
\begin{itemize}[leftmargin=2em]
    \item coverage at nominal levels such as $50\%$, $80\%$, $90\%$, and $95\%$;
    \item interval width, which controls sharpness;
    \item Gaussian NLL, which rewards both accuracy and calibrated uncertainty;
    \item probability integral transform (PIT) diagnostics for calibration shape.
\end{itemize}

If $F_j$ is the predictive CDF for observation $y_j$, then
\[
    \PIT_j = F_j(y_j)
\]
should be approximately uniform on $[0,1]$ under perfect calibration.

\subsection{Why Notebook 03 Matters}

Notebook 03 is theoretically important because it shifts the goal from
\[
    \text{best deterministic forecast}
\]
to
\[
    \text{best calibrated predictive distribution}.
\]
That shift becomes even more important once later notebooks show that raw mechanistic parameters are weakly identifiable. Even when parameters are poorly recovered, useful predictive uncertainty may still be achievable.

Notebook 04 operationalizes these ideas at full cohort scale and under leakage-free evaluation.

% -----------------------------------------------------------------------------
\section{Notebook 04: Cohort-Scale Leakage-Free Baselines and SBI Readiness}
% -----------------------------------------------------------------------------

Notebook 04 scales the rolling-origin and calibration framework from a 20-patient exploratory cohort to the canonical cross-study dataset.

\subsection{Cohort Definition}

The forecasting cohort applies stricter filters:
\[
    \text{at least 6 scans},\qquad
    \text{duration at least 90 days},\qquad
    \text{minimum training scans }=4.
\]
This produced:
\begin{itemize}[leftmargin=2em]
    \item 5,862 patients in the overall feature table,
    \item 1,923 patients in the forecasting cohort,
    \item 7,710 deterministic one-step predictions,
    \item 23,130 long-format prediction rows across multiple models and calibration views.
\end{itemize}

\subsection{Leakage-Free Calibration Pools}

Residual pools are formed without using the target patient's own future prediction error. Notebook 04 tests several calibration scopes:
\begin{itemize}[leftmargin=2em]
    \item \texttt{global\_lopo}: global leave-one-patient-out,
    \item \texttt{gapbin\_global\_lopo}: global but stratified by scan-gap bin,
    \item \texttt{study\_lopo}: leave-one-patient-out within study,
    \item \texttt{drug\_lopo}: leave-one-patient-out within drug signature,
    \item \texttt{study\_drug\_lopo}: leave-one-patient-out within study-drug grouping.
\end{itemize}

The goal is not only calibration accuracy but also realism: uncertainty should adapt to the context in which the forecast is made.

\subsection{Deterministic Baselines at Scale}

The cohort-scale deterministic results are shown in Table~\ref{tab:n04-det}.

\begin{table}[H]
\centering
\small
\caption{Notebook 04 deterministic empirical baselines over the cohort-scale rolling-origin task.}
\label{tab:n04-det}
\begin{tabular}{lrrrr}
\toprule
Model & Predictions & Patients & Median $|\log$ err$|$ & RMSE \\
\midrule
LOCF & 7,710 & 1,923 & 0.068 & 0.302 \\
RecentLogSlope & 7,710 & 1,923 & 0.105 & 0.359 \\
LogLinear & 7,710 & 1,923 & 0.224 & 0.438 \\
\bottomrule
\end{tabular}
\end{table}

LOCF remains the best deterministic model. This is an important continuity check with Notebook 02: the dominance of persistence was not a small-cohort fluke.

\subsection{Probabilistic Calibration Outputs}

Notebook 04 identifies two especially relevant calibrated baselines.

\paragraph{Best Gaussian calibration.}
The best Gaussian calibration uses LOCF with \texttt{gapbin\_global\_lopo}. Its key values are:
\[
\begin{aligned}
    \text{RMSE} &= 0.3015,\\
    \text{mean Gaussian NLL} &= 0.2032,\\
    \text{coverage}_{80} &= 0.9043,\\
    \text{coverage}_{90} &= 0.9333,\\
    \text{coverage}_{95} &= 0.9514.
\end{aligned}
\]
The 95\% interval is almost perfectly calibrated, while lower nominal intervals are conservative.

\paragraph{Best empirical-quantile calibration.}
The most accurate nominal coverage comes from RecentLogSlope with \texttt{global\_lopo} empirical quantiles:
\[
\begin{aligned}
    \text{coverage}_{50} &= 0.4997,\\
    \text{coverage}_{80} &= 0.7997,\\
    \text{coverage}_{90} &= 0.8999,\\
    \text{coverage}_{95} &= 0.9499,
\end{aligned}
\]
with mean absolute coverage error essentially zero. This calibration is extremely well matched to nominal probabilities, though the underlying deterministic forecaster is weaker than LOCF.

Table~\ref{tab:n04-calib} summarizes the two main baselines.

\begin{table}[H]
\centering
\small
\caption{Notebook 04 strongest leakage-free calibrated baselines.}
\label{tab:n04-calib}
\begin{tabular}{lrrrrrrrr}
\toprule
Model & Scope & Interval & Cov.50 & Cov.80 & Cov.90 & Cov.95 & Width80 & Mean cov. err \\
\midrule
LOCF & gapbin\_global\_lopo & Gaussian & 0.781 & 0.904 & 0.933 & 0.951 & 0.654 & 0.105 \\
RecentLogSlope & global\_lopo & Emp.\ quantile & 0.500 & 0.800 & 0.900 & 0.950 & 0.520 & 0.00018 \\
\bottomrule
\end{tabular}
\end{table}

\subsection{SBI Readiness by Study/Drug/Arm Stratum}

Notebook 04 does more than score forecasting baselines. It introduces a concrete SBI-readiness layer. The rationale is:
\[
\text{Do not train SBI across all trials indiscriminately.}
\]
Instead, select coherent study/drug/arm strata with adequate scan depth, dose availability, manageable baseline error, and interpretable mechanism hints.

The readiness counts were:
\begin{itemize}[leftmargin=2em]
    \item green: 8 strata,
    \item yellow: 8 strata,
    \item red: 17 strata.
\end{itemize}

The top recommended starter strata are shown in Table~\ref{tab:n04-readiness}.

\begin{table}[H]
\centering
\small
\caption{Top SBI starter strata from Notebook 04.}
\label{tab:n04-readiness}
\begin{tabular}{p{0.18\textwidth}p{0.18\textwidth}p{0.16\textwidth}rrrr}
\toprule
Study & Arm & Family hint & Patients & Median scans & Baseline RMSE & Cov.90 \\
\midrule
CEDIRANIB\_HORIZONIII & FOLFOX + bevacizumab 5mg/kg & anti-angiogenic & 268 & 6 & 0.297 & 0.931 \\
DATASET\_264 & FOLFOX alone & cytotoxic TGI & 190 & 7 & 0.361 & 0.891 \\
EFC10262 & Placebo + FOLFIRI & cytotoxic TGI & 190 & 8 & 0.358 & 0.932 \\
DATASET\_264 & Panitumumab + FOLFOX & cytotoxic TGI & 188 & 8 & 0.374 & 0.867 \\
DATASET\_263 & Panitumumab + FOLFIRI & cytotoxic TGI & 152 & 7 & 0.306 & 0.915 \\
A6181122 & UNKNOWN\_ARM & cytotoxic TGI & 127 & 8 & 0.331 & 0.911 \\
EFC5505 & PLACEBO & cytotoxic TGI & 120 & 7 & 0.498 & 0.877 \\
DATASET\_263 & FOLFIRI alone & cytotoxic TGI & 104 & 7 & 0.335 & 0.907 \\
\bottomrule
\end{tabular}
\end{table}

\subsection{Why Notebook 04 Is a Turning Point}

Notebook 04 formalizes three project-level rules:
\begin{enumerate}[leftmargin=2em]
    \item the official predictive baseline is leakage-free calibrated empirical forecasting, not retrospective ODE fit;
    \item SBI should begin in coherent green/yellow strata, not across the entire pooled dataset;
    \item SBI should be evaluated by posterior predictive calibration, NLL, interval width, and long-horizon behavior.
\end{enumerate}

The top green stratum was
\[
    \texttt{CEDIRANIB\_HORIZONIII / FOLFOX + bevacizumab 5 mg/kg},
\]
which becomes the anchor of Notebook 05.

% -----------------------------------------------------------------------------
\section{Literature and Model Catalogue}
% -----------------------------------------------------------------------------

The literature/model catalogue runs in parallel with the forecasting notebooks. Its role is to stop the project from adding compartments blindly just because a richer model helped one patient.

\subsection{Mechanism-to-Model Routing}

The catalogue maps drug classes and observed trajectory patterns to model families:
\begin{itemize}[leftmargin=2em]
    \item cytotoxic chemotherapy $\rightarrow$ Simeoni TGI or nested resistant-Simeoni,
    \item targeted therapy $\rightarrow$ Emax or resistance-aware TGI,
    \item anti-angiogenic therapy $\rightarrow$ carrying-capacity / vascular-support models,
    \item immune checkpoint therapy $\rightarrow$ delayed immune-response models,
    \item combination therapy $\rightarrow$ multi-input interaction TGI,
    \item lesion-rich data $\rightarrow$ lesion-level hierarchical models.
\end{itemize}

\subsection{Why Anti-Angiogenic Therapy Needs a Different Latent Story}

The catalogue explicitly warns against representing anti-VEGF-like treatment as direct cytotoxic kill alone. A reduced carrying-capacity model uses tumor burden $V(t)$ and vascular support $K(t)$:
\[
\begin{aligned}
    \frac{dV}{dt}&=\rho V\left(1-\frac{V}{K}\right),\\
    \frac{dK}{dt}&=k_{K,\mathrm{recovery}}(K_0-K)-k_{\mathrm{anti}}C(t)K.
\end{aligned}
\]
This is the conceptual starting point for the anti-angiogenic simulator later used in Notebooks 05C and 06.

\subsection{Strategic Use of the Catalogue}

The catalogue changes model development from
\[
    \text{``what extra compartment should we add next?''}
\]
to
\[
    \text{``what latent structure is justified by drug mechanism, data granularity, and identifiability?''}
\]
That is why Notebook 05 does not simply export the original all-purpose Simeoni model. It exports simulator families aligned with specific strata.

% -----------------------------------------------------------------------------
\section{Notebook 05: Selected-Stratum Simulator Design and Synthetic SBI Preparation}
% -----------------------------------------------------------------------------

Notebook 05 is where the project transitions from real-data forecasting diagnostics to synthetic SBI design.

\subsection{Selected Target Stratum}

The chosen first stratum is
\[
    \texttt{CEDIRANIB\_HORIZONIII / FOLFOX + bevacizumab 5 mg/kg}.
\]
Its readiness summary from Notebook 04 is:
\begin{itemize}[leftmargin=2em]
    \item 268 real patients,
    \item median 6 scans,
    \item median duration 362 days,
    \item median gap 64 days,
    \item full dose availability,
    \item baseline RMSE 0.297,
    \item baseline 90\% coverage 0.931.
\end{itemize}

The response-class fractions in the real cohort are:
\[
\begin{aligned}
    \text{response or shrinkage} &= 0.459,\\
    \text{response then regrowth} &= 0.325,\\
    \text{stable-like} &= 0.198,\\
    \text{progression or growth} &= 0.019.
\end{aligned}
\]

\subsection{Synthetic Simulator Philosophy}

Notebook 05 does not simulate artificial evenly spaced trajectories from generic priors alone. It reuses real patient templates:
\begin{itemize}[leftmargin=2em]
    \item real scan schedules,
    \item real dose schedules,
    \item real trajectory lengths,
    \item real baseline scales,
    \item clinically coherent study/drug/arm metadata.
\end{itemize}

This is crucial. For SBI, the observation design itself is part of the inverse problem.

\subsection{Simulator Families}

The first 05 export used two simulator families.

\paragraph{Simeoni family.}
This retains the cytotoxic transit-damage model and serves as a mechanistic baseline family.

\paragraph{Angio family.}
For anti-angiogenic behavior, a reduced tumor-plus-carrying-capacity model is used:
\[
\begin{aligned}
    \frac{dV}{dt}
    &=\rho V\left(1-\frac{V}{K}\right),\\
    \frac{dK}{dt}
    &=k_{K,\mathrm{recovery}}(K_0-K)-k_{\mathrm{anti}}C(t)K.
\end{aligned}
\]
The initial capacity is parameterized as
\[
    K(0)=K_0=\text{K0\_ratio}\cdot V(0),
\]
and the core angio parameter vector becomes
\[
    \theta_{\mathrm{angio}}
    =
    (\rho,k_{\mathrm{anti}},k_{K,\mathrm{recovery}},K0\_ratio,v_0).
\]

This reduced model is intentionally low-dimensional. It is designed to test whether anti-angiogenic local dynamics can be inferred at all from sparse relative \sld{} trajectories.

\subsection{Initial 05 Export}

The first synthetic export contained:
\begin{itemize}[leftmargin=2em]
    \item 2,000 total simulations,
    \item 1,000 Simeoni simulations,
    \item 1,000 angio simulations,
    \item array shape $\xobs\in\R^{2000\times 14\times 3}$,
    \item parameter array shape $\theta\in\R^{2000\times 10}$ in the mixed-family export.
\end{itemize}

The mixed theta names were
\[
(\theta_{K0\_ratio},\theta_{k_K\_recovery},\theta_{k\_anti},\theta_{k\_cyt},\theta_{k\_tr},\theta_{\lambda0},\theta_{\lambda1},\theta_{\psi_g},\theta_{\rho},\theta_{v0}),
\]
with a simulator ID indicating whether the sample came from the Simeoni or angio family.

\subsection{Why 05 Needed Further Tuning}

The first prior-predictive cohort did not match the real stratum closely enough. In the target CEDIRANIB stratum, the synthetic response-class fractions were too progression-heavy:
\begin{itemize}[leftmargin=2em]
    \item angio progression/growth: $0.422$,
    \item Simeoni progression/growth: $0.557$,
    \item real progression/growth: only $0.0187$.
\end{itemize}

That mismatch is not a cosmetic issue. If the synthetic generator produces the wrong local dynamics, an NPE trained on it will learn the wrong inverse mapping.

% -----------------------------------------------------------------------------
\section{Notebook 05B and 05C: Tuned Multistratum Synthetic Exports}
% -----------------------------------------------------------------------------

The 05B and 05C updates are where the simulator preparation becomes mature enough for Notebook 06.

\subsection{05B: Prior Tuning and Multistratum Export}

Notebook 05B expands from one target stratum to four selected strata and tunes class quotas so the synthetic response-pattern mixture better resembles the real data.

The four strata are:
\begin{enumerate}[leftmargin=2em]
    \item CEDIRANIB\_HORIZONIII / FOLFOX + bevacizumab 5 mg/kg,
    \item DATASET\_264 / FOLFOX alone,
    \item EFC10262 / Placebo + FOLFIRI,
    \item DATASET\_264 / Panitumumab + FOLFOX.
\end{enumerate}

The tuned export contains
\[
    4{,}800
\]
synthetic trajectories total:
\[
    2{,}400\ \text{Simeoni}
    \qquad\text{and}\qquad
    2{,}400\ \text{angio}.
\]
Each study-arm-stratum contributes 600 simulations per simulator family.

The grouped split is by real template rather than by individual simulation row, which is essential to prevent template leakage between train and test.

\subsection{05C: Local-Dynamics Matching}

Notebook 05C refines the export again, this time emphasizing local one-step dynamics and real-vs-synthetic residual matching. The final 05C exports are:
\[
    \xobs\in\R^{2400\times 25\times 3},
\]
\[
    x_{\mathrm{cov}}\in\R^{2400\times 8},
\]
with covariates
\[
    (\log \text{baseline \sld},\ \text{duration\_years},\ n\_scans/25,\ \text{dose\_confounded},\ \text{4 stratum one-hots}).
\]

The train/validation/test split sizes for the angio export are:
\[
    1948,\quad 218,\quad 234.
\]

\subsection{05C Real-vs-Synthetic Output Matching}

The response summaries in Table~\ref{tab:05c-response} show how much closer the tuned synthetic cohorts came to the real strata.

\begin{table}[H]
\centering
\small
\caption{Selected 05C real-vs-synthetic response summaries.}
\label{tab:05c-response}
\begin{tabular}{p{0.21\textwidth}p{0.10\textwidth}rrrr}
\toprule
Source / stratum & Simulator & Median final rel.\ \sld{} & Median minimum rel.\ \sld{} & Shrink frac & Regrowth frac \\
\midrule
CEDIRANIB real & real & 0.629 & 0.521 & 0.459 & 0.325 \\
CEDIRANIB 05C & angio & 0.459 & 0.392 & 0.453 & 0.322 \\
CEDIRANIB 05C & simeoni & 0.311 & 0.256 & 0.453 & 0.322 \\
FOLFOX-alone real & real & 0.619 & 0.404 & 0.389 & 0.458 \\
FOLFOX-alone 05C & angio & 0.525 & 0.381 & 0.390 & 0.458 \\
FOLFOX-alone 05C & simeoni & 0.286 & 0.212 & 0.390 & 0.458 \\
\bottomrule
\end{tabular}
\end{table}

The local residual behavior is even more informative. For the target CEDIRANIB stratum:
\[
\text{real LOCF RMSE}=0.231,
\qquad
\text{angio synthetic LOCF RMSE}=0.176,
\qquad
\text{simeoni synthetic LOCF RMSE}=0.486.
\]
Thus the angio simulator is much closer than the Simeoni family to the local dynamics of the anti-angiogenic real cohort.

Table~\ref{tab:05c-residual} summarizes the key values.

\begin{table}[H]
\centering
\small
\caption{05C local-dynamics matching via LOCF residual summaries.}
\label{tab:05c-residual}
\begin{tabular}{p{0.21\textwidth}p{0.10\textwidth}rrrr}
\toprule
Source / stratum & Simulator & Count & Median $|\log$ err$|$ & RMSE & $q_{95}$ \\
\midrule
CEDIRANIB real & real & 772 & 0.064 & 0.231 & 0.363 \\
CEDIRANIB synthetic & angio & 1,670 & 0.112 & 0.176 & 0.333 \\
CEDIRANIB synthetic & simeoni & 1,661 & 0.250 & 0.486 & 0.555 \\
EFC10262 real & real & 850 & 0.043 & 0.384 & 0.338 \\
EFC10262 synthetic & angio & 2,306 & 0.077 & 0.144 & 0.262 \\
EFC10262 synthetic & simeoni & 2,162 & 0.131 & 0.307 & 0.349 \\
\bottomrule
\end{tabular}
\end{table}

\subsection{Why 05C Becomes the Official Notebook 06 Input}

Notebook 05C is the correct launch point for NPE because it satisfies three conditions:
\begin{enumerate}[leftmargin=2em]
    \item grouped template splits avoid train/test leakage,
    \item local dynamics are substantially closer to real data than the naive export,
    \item the angio simulator matches the anti-angiogenic target stratum far better than the plain Simeoni family.
\end{enumerate}

That is why Notebook 06 uses the 05C angio-only export rather than the earlier mixed-family array.

% -----------------------------------------------------------------------------
\section{Notebook 06: First Neural Posterior Estimation on 05C Angio}
% -----------------------------------------------------------------------------

Notebook 06 asks the first truly SBI-specific question:
\[
    \text{Can a learned posterior recover simulator parameters from synthetic tumor trajectories?}
\]

\subsection{Input Representation}

For each simulation, the padded observation sequence is
\[
    \xobs
    =
    [\tilde t,\log r,m]
    \in\R^{T_{\max}\times 3},
\]
where $\tilde t$ is normalized time, $\log r$ is log-relative tumor burden, and $m$ is a mask for variable-length trajectories.

The covariates are
\[
    x_{\mathrm{cov}}\in\R^8,
\]
with baseline scale, duration, scan count, dose-confounding flag, and one-hot stratum indicators.

The 05C angio target parameters are
\[
    \theta=(\rho,k_{\mathrm{anti}},k_{K,\mathrm{recovery}},K0\_ratio,v_0).
\]
Notebook 06 standardizes them on log scale:
\[
    z=\frac{\log\theta-\mu_{\log\theta}}{\sigma_{\log\theta}}.
\]

\subsection{Posterior Families and Baselines}

Notebook 06 compares:
\begin{itemize}[leftmargin=2em]
    \item a prior Gaussian baseline,
    \item Ridge residual bootstrap,
    \item Random Forest residual bootstrap,
    \item summary-feature MLP regression,
    \item summary-feature MDNs,
    \item GRU-based MDNs.
\end{itemize}

The core posterior density family is a mixture density network:
\[
    q_\phi(z\given x)
    =
    \sum_{k=1}^{K}\pi_k(x)\Normal\!\bigl(z;\mu_k(x),\diag(\sigma_k^2(x))\bigr).
\]
The GRU variant uses the sequential trajectory directly; the summary models compress the trajectory into engineered features first.

\subsection{Train/Validation/Test Split}

The 05C angio export used here contains:
\begin{itemize}[leftmargin=2em]
    \item $N=2400$ total simulations,
    \item $T_{\max}=25$,
    \item train/val/test sizes $1948/218/234$,
    \item 757 unique real templates,
    \item four strata.
\end{itemize}

Because the split is template-grouped, the task is harder and more realistic than a random row split.

\subsection{Notebook 06 Aggregate Results}

The selected best model is
\[
    \texttt{gru\_mdn\_K5\_H128}.
\]
Its aggregate held-out metrics are
\[
\begin{aligned}
    \text{mean RMSE}_z &= 0.858,\\
    \text{mean MAE}_z &= 0.727,\\
    \text{mean RMSE}_{\log\theta} &= 0.883,\\
    \text{mean abs.\ relative error} &= 0.981,\\
    \text{mean corr}_z &= 0.454,\\
    \text{mean coverage}_{90} &= 0.915,\\
    \text{mean posterior std}_z &= 0.866.
\end{aligned}
\]

This is better than the prior and slightly better than the Ridge / Random Forest baselines, but the margin is not large.

Parameter-level performance is shown in Table~\ref{tab:n06-param}.

\begin{table}[H]
\centering
\small
\caption{Notebook 06 best-model parameter recovery on held-out 05C angio simulations.}
\label{tab:n06-param}
\begin{tabular}{lrrrrr}
\toprule
Parameter & RMSE$_z$ & Corr$_z$ & Mean rel.\ err & Cov.90 & Posterior std$_z$ \\
\midrule
$\rho$ & 0.905 & 0.450 & 0.928 & 0.906 & 0.913 \\
$k_{\mathrm{anti}}$ & 0.928 & 0.384 & 1.506 & 0.906 & 0.901 \\
$k_{K,\mathrm{recovery}}$ & 0.908 & 0.377 & 1.315 & 0.944 & 0.939 \\
$K0\_ratio$ & 0.952 & 0.265 & 0.386 & 0.915 & 0.950 \\
$v_0$ & 0.598 & 0.792 & 0.770 & 0.902 & 0.625 \\
\bottomrule
\end{tabular}
\end{table}

\subsection{Interpretation of Notebook 06}

Only $v_0$ is recovered moderately well, with correlation near
\[
    0.79.
\]
This is not surprising because baseline tumor scale is directly represented in the covariates. The harder mechanistic parameters
\[
    \rho,\quad k_{\mathrm{anti}},\quad k_{K,\mathrm{recovery}},\quad K0\_ratio
\]
are recovered only weakly. Correlations are low to moderate and errors remain substantial.

This result is extremely important:
\[
\boxed{
\text{even under simulator-matched synthetic data, the angio parameters are only weakly recoverable}
}
\]
Therefore the main problem is not just neural architecture quality. It is the information content of sparse relative \sld{} trajectories.

% -----------------------------------------------------------------------------
\section{Notebook 06B: Stabilized NPE and Identifiability Rescue Attempts}
% -----------------------------------------------------------------------------

Notebook 06B is explicitly a diagnosis notebook. It asks whether poor performance is due to:
\begin{enumerate}[leftmargin=2em]
    \item random-vs-grouped split shift,
    \item unstable MDN training,
    \item excessive target dimensionality,
    \item bad parameterization of the latent dynamics,
    \item or fundamental weak identifiability.
\end{enumerate}

\subsection{Stabilization Strategy}

The patch replaces unstable MDN heads with simpler Gaussian posterior heads and tests reduced/composite targets. Instead of always predicting the full five raw angio parameters, Notebook 06B defines:
\begin{itemize}[leftmargin=2em]
    \item \texttt{core2\_growth\_treatment}: $(\log\rho,\log k_{\mathrm{anti}})$,
    \item \texttt{core3\_logtheta}: $(\log\rho,\log k_{\mathrm{anti}},\log K0\_ratio)$,
    \item \texttt{composite3}: growth, treatment-vs-growth, recovery-vs-treatment, and capacity coordinates.
\end{itemize}

The composite representation is meant to test whether the data identify \emph{relative dynamical balances} better than raw rate constants.

\subsection{Decision Rule}

The patch compares each stabilized model to a prior baseline on the same target. If the relative RMSE improvement over the prior is small, the model is still weakly informative.

In general form, define
\[
    G
    =
    \frac{\mathrm{RMSE}_{\mathrm{prior}}-\mathrm{RMSE}_{\mathrm{model}}}
    {\mathrm{RMSE}_{\mathrm{prior}}}.
\]
Notebook 06B interprets
\[
    5\%\text{--}15\%
\]
relative gain as a weak-identifiability regime rather than a strong inference success.

\subsection{Notebook 06B Outputs}

The strongest reduced/composite targets are summarized in Table~\ref{tab:n06b}.

\begin{table}[H]
\centering
\small
\caption{Notebook 06B stabilized Gaussian target-comparison summary.}
\label{tab:n06b}
\begin{tabular}{p{0.21\textwidth}rrrccc}
\toprule
Target & Dim & Best model & RMSE$_z$ & Prior RMSE$_z$ & Rel.\ gain & Cov.90 \\
\midrule
composite3 & 4 & GRU Gaussian & 0.900 & 1.006 & 10.5\% & 0.868 \\
core2\_growth\_treatment & 2 & Summary Gaussian & 0.907 & 1.007 & 9.95\% & 0.853 \\
core3\_logtheta & 3 & GRU Gaussian & 0.924 & 1.001 & 7.72\% & 0.886 \\
\bottomrule
\end{tabular}
\end{table}

The best configuration is
\[
    \texttt{06B\_gru\_gaussian\_composite\_global\_sigma},
\]
with
\[
    \text{mean RMSE}_z=0.900,
    \qquad
    \text{relative improvement over prior}=10.5\%.
\]

However, that still falls into the weak-identifiability band defined by the notebook.

\subsection{Interpretation of Notebook 06B}

Notebook 06B is one of the most important notebooks in the whole sequence because it prevents the wrong conclusion. A naive reading of Notebook 06 might say:
\[
    \text{``try a larger network''}.
\]
Notebook 06B shows why that is the wrong move. Reduced targets, composite targets, and stabilized Gaussian heads only produce modest improvements over the prior:
\[
    7.7\%\text{ to }10.5\%.
\]
That is a scientific signal, not a training inconvenience. It means:
\[
\boxed{
\text{sparse relative \sld{} trajectories weakly identify composite growth/treatment balances, but not detailed raw angio parameters}
}
\]

The right next step is therefore posterior predictive validation, not parameter-regression escalation.

% -----------------------------------------------------------------------------
\section{Integrated Scientific Interpretation}
% -----------------------------------------------------------------------------

Across Notebooks 00--06B, the project asks increasingly difficult questions.

\begin{enumerate}[leftmargin=2em]
    \item \textbf{Can the raw clinical files be made analysis-ready?}
    Yes. Ten studies were canonicalized, with 5,862 patients entering the feature table.

    \item \textbf{Can mechanistic models fit retrospective trajectories?}
    Yes. On a selected 20-patient batch, median retrospective RMSE improved from $0.290$ under plain Simeoni to $0.151$ after conservative selection over nested resistance and observation shifts.

    \item \textbf{Do those fits generalize to future scans?}
    Only weakly. Hold-out median test RMSEs remained around $0.49$--$0.51$.

    \item \textbf{Do mechanistic models beat simple forecasting baselines?}
    No. LOCF dominated rolling-origin point forecasting on both the 20-patient and 1,923-patient cohorts.

    \item \textbf{Can useful predictive uncertainty be calibrated empirically?}
    Yes. Leakage-free Gaussian and empirical-quantile calibration produced strong coverage baselines.

    \item \textbf{Can a clinically coherent SBI starter problem be isolated?}
    Yes. The CEDIRANIB anti-angiogenic stratum was selected as the first green starter stratum.

    \item \textbf{Can synthetic training data be made structurally similar to the real local dynamics?}
    Partly yes. The 05C angio export is much closer than the naive or Simeoni-only synthetic alternatives.

    \item \textbf{Can NPE recover raw simulator parameters on simulator-matched synthetic data?}
    Only weakly. The best GRU-MDN barely outperformed simpler baselines, and 06B showed only modest gains over the prior even for reduced targets.
\end{enumerate}

The full project therefore delivers a valuable negative result:
\[
\boxed{
\text{the real difficulty is not just flexible modeling; it is weak identifiability from sparse tumor trajectories}
}
\]

That negative result is productive. It stops the project from over-interpreting neural posterior outputs and redirects it toward better scientific evaluation criteria.

% -----------------------------------------------------------------------------
\section{What the Next Full Notebook Should Be}
% -----------------------------------------------------------------------------

The next notebook should not be another architecture search. It should be a posterior predictive notebook:
\[
    \texttt{07\_angio\_posterior\_predictive\_checks\_synthetic\_then\_real.ipynb}.
\]

\subsection{Goal}

The target question is
\[
    \text{Can posterior samples generate calibrated predictive bands over tumor trajectories?}
\]
not
\[
    \text{Can we recover every mechanistic parameter precisely?}
\]

\subsection{Posterior Predictive Workflow}

For a fitted posterior $q_\phi(\theta\given x)$, the posterior predictive procedure is
\[
    \theta^{(s)}\sim q_\phi(\theta\given x),
\]
\[
    \widetilde y^{(s)}_{1:T}\sim p(y_{1:T}\given \theta^{(s)},u),
\]
followed by comparisons between the empirical trajectory and the posterior predictive ensemble.

\subsection{Pass/Fail Metrics}

The next notebook should report:
\begin{itemize}[leftmargin=2em]
    \item predictive coverage at $80\%$, $90\%$, and $95\%$,
    \item interval width versus the empirical baselines from Notebook 04,
    \item posterior predictive NLL,
    \item calibration by stratum and scan-gap bin,
    \item robustness to sparse prefixes and local jumps,
    \item synthetic-first validation before any real-patient interpretation.
\end{itemize}

The standard for success is not ``posterior samples look diverse.'' The standard is:
\[
\text{posterior predictive performance at least matches calibrated empirical baselines}
\]
and ideally improves on them in clinically coherent strata.

% -----------------------------------------------------------------------------
\section{Limitations}
% -----------------------------------------------------------------------------

The current pipeline remains a deliberately simplified but rigorous pre-SBI program. Main limitations are:
\begin{itemize}[leftmargin=2em]
    \item the exposure model is still an exponential forcing approximation, not full PK;
    \item \sld{} compresses lesion-level heterogeneity into a single scalar;
    \item observation shifts are phenomenological and may reflect multiple non-biological processes;
    \item synthetic angio dynamics are reduced effective dynamics, not a full vascular biology model;
    \item parameter inference is tested mainly on synthetic data and remains weak even there;
    \item the strongest real-data predictive baselines remain empirical rather than mechanistic.
\end{itemize}

These are not failures. They define the correct next stage of work.

% -----------------------------------------------------------------------------
\section{Conclusion}
% -----------------------------------------------------------------------------

The draft \texttt{notebook1\_all\_notebooks\_updated.tex} has now been expanded into a full report that treats each notebook as a real modeling stage rather than a short summary bullet. The complete notebook sequence tells a coherent story:

\begin{enumerate}[leftmargin=2em]
    \item clinical data can be canonicalized across heterogeneous oncology trials;
    \item mechanistic TGI models can fit many retrospective trajectories well;
    \item those fits do not guarantee good future prediction;
    \item strong empirical rolling-origin baselines are hard to beat;
    \item calibrated predictive uncertainty is feasible and should be the benchmark;
    \item simulator choice must follow drug mechanism and data context;
    \item synthetic SBI is feasible in carefully selected strata;
    \item raw parameter recovery remains weak even under simulator-matched synthetic data.
\end{enumerate}

The most important conclusion is therefore:
\[
\boxed{
\begin{minipage}{0.86\textwidth}
Sparse relative \sld{} trajectories are informative enough to support useful empirical forecasting and some coarse dynamical uncertainty quantification, but they do not strongly identify detailed mechanistic PK/PD parameters. The next scientifically meaningful milestone is posterior predictive calibration, not aggressive raw-parameter recovery.
\end{minipage}
}
\]

That is a strong methodological result. It protects future SBI work from over-claiming and defines a realistic path toward a credible oncology posterior simulator.

\appendix

% -----------------------------------------------------------------------------
\section{Notebook-to-Output Map}
% -----------------------------------------------------------------------------

\begin{center}
\small
\begin{tabular}{p{0.16\textwidth}p{0.30\textwidth}p{0.36\textwidth}}
\toprule
Notebook & Main purpose & Most important saved output or conclusion \\
\midrule
00 & Clinical audit and first mechanistic discovery & Ten studies canonicalized; retrospective 20-patient median RMSE improved to $0.151$ after conservative model selection. \\
01 & Hold-out validation & Future-scan median test RMSEs remained near $0.49$--$0.51$, exposing overfit risk. \\
02 & Rolling-origin forecasting & LOCF dominated with RMSE $0.185$ on the 20-patient cohort. \\
03 & Probabilistic calibration theory & Formalized Gaussian and empirical residual calibration, coverage, NLL, and PIT. \\
04 & Cohort-scale leakage-free baselines & 1,923-patient cohort; LOCF best deterministic baseline; near-perfect empirical interval coverage available. \\
05 & Selected-stratum simulator design & CEDIRANIB anti-angiogenic starter stratum selected; first synthetic SBI arrays exported. \\
05B & Multistratum tuned export & Four-stratum, grouped-split synthetic export with 2,400 angio and 2,400 Simeoni simulations. \\
05C & Local-dynamics matching & Angio simulator brought much closer to real local residual behavior than Simeoni in the target anti-angiogenic stratum. \\
06 & First angio NPE & Best GRU-MDN achieved mean RMSE$_z=0.858$ but weak raw parameter recovery overall. \\
06B & Identifiability rescue attempt & Reduced/composite targets improved only $7.7\%$--$10.5\%$ over the prior, implying weak identifiability. \\
\bottomrule
\end{tabular}
\end{center}

% -----------------------------------------------------------------------------
\begin{thebibliography}{12}

\bibitem{simeoni2004}
M. Simeoni et al., ``Predictive pharmacokinetic-pharmacodynamic modeling of tumor growth kinetics in xenograft models after administrations of anticancer agents,'' \emph{Cancer Research}, 2004.

\bibitem{magni2008}
P. Magni et al., ``A minimal model of tumor growth inhibition,'' 2008.

\bibitem{terranova2013}
N. Terranova et al., ``A predictive pharmacokinetic-pharmacodynamic model of tumor growth kinetics in xenograft mice after administration of anticancer agents given in combination,'' 2013.

\bibitem{claret2009}
L. Claret et al., ``Model-based prediction of phase III overall survival in colorectal cancer on the basis of phase II tumor dynamics,'' \emph{Journal of Clinical Oncology}, 2009.

\bibitem{krishnan2021}
S. Krishnan et al., ``Tumor growth inhibition modeling of individual lesion dynamics and inter-organ variability in HER2-negative breast cancer patients treated with docetaxel,'' 2021.

\bibitem{courlet2023}
P. Courlet et al., ``Modeling tumor size dynamics based on real-world clinical and imaging data in advanced melanoma patients receiving immune checkpoint inhibitors,'' 2023.

\bibitem{oden2013}
J. T. Oden et al., ``Selection and assessment of phenomenological models of tumor growth,'' 2013.

\bibitem{whitaker2019}
G. A. Whitaker et al., ``Bayesian inference for stochastic differential equation mixed effects models of a tumour xenography study,'' 2019.

\bibitem{ezhov2023}
D. Ezhov et al., ``Learn-Morph-Infer: A new way of solving the inverse problem for brain tumor modeling,'' 2023.

\bibitem{yankeelov2013}
T. E. Yankeelov et al., ``Clinically Relevant Modeling of Tumor Growth and Treatment Response,'' \emph{Science Translational Medicine}, 2013.

\bibitem{gholami2016}
A. Gholami et al., ``An inverse problem formulation for parameter estimation of a reaction-diffusion model of low grade gliomas,'' \emph{Journal of Mathematical Biology}, 2016.

\bibitem{le2018}
A. L\^e et al., ``Subject-specific brain tumor growth modelling via an efficient Bayesian inference framework,'' in \emph{Medical Imaging 2018: Image Processing}, 2018.

\end{thebibliography}

\end{document}
